\documentclass[11pt,a4paper]{article}
\usepackage[utf8]{inputenc}
\usepackage[T1]{fontenc}
\usepackage[spanish]{babel}
\usepackage{libertinus}
\usepackage{microtype}
\usepackage[margin=1.8cm]{geometry}
\usepackage{graphicx}
\usepackage{xcolor}
\usepackage[hidelinks]{hyperref}
\usepackage{titlesec}
\definecolor{UGblue}{RGB}{0,51,102}
\definecolor{UGgold}{RGB}{212,175,55}
\titleformat{\section}{\Large\bfseries\color{UGblue}}{}{0pt}{}[{\color{UGgold}\titlerule[1pt]}]
\setlength{\parindent}{0pt}
% Cada figura: miniatura escalada que cabe en una celda de la cuadrícula + nombre del archivo.
\newcommand{\fig}[2]{%
  \begin{minipage}[b]{0.31\textwidth}\centering
    \includegraphics[width=\linewidth,height=4.2cm,keepaspectratio]{#1/#2.pdf}\\[2pt]
    {\footnotesize\ttfamily\color{black!60}#2}
  \end{minipage}\hfill\ignorespaces}
\begin{document}
\begin{titlepage}\centering
  \vspace*{5cm}
  {\Huge\bfseries\color{UGblue} Catálogo de figuras\par}
  \vspace{0.4cm}{\color{UGgold}\rule{0.5\textwidth}{1.2pt}}\par\vspace{0.4cm}
  {\Large TikZ · PGFPlots · tikz-cd\par}
  \vspace{1.5cm}
  {\large Ricardo León Martínez\par}
  {\normalsize Departamento de Matemáticas · Universidad de Guanajuato\par}
  \vfill
  {\small 215 figuras generadas desde código\par}
\end{titlepage}
\tableofcontents
\newpage

\section{Cálculo I}
{\small\color{black!60}81 figuras · \texttt{galeria/calculo-I/}}\par\bigskip

\fig{calculo-I}{analisis-quintica}
\fig{calculo-I}{asintota-horizontal}
\fig{calculo-I}{asintota-oblicua}
\par\bigskip
\fig{calculo-I}{asintota-vertical}
\fig{calculo-I}{caja}
\fig{calculo-I}{circunferencia}
\par\bigskip
\fig{calculo-I}{concavidad-opciones}
\fig{calculo-I}{criterio-primera}
\fig{calculo-I}{cuartica-w}
\par\bigskip
\fig{calculo-I}{cubica-crece}
\fig{calculo-I}{cubica-inflexion}
\fig{calculo-I}{curva-asintotica}
\par\bigskip
\fig{calculo-I}{curva-gamma}
\fig{calculo-I}{derivada-fisica}
\fig{calculo-I}{derivada-geometrica}
\par\bigskip
\fig{calculo-I}{elipse-tangente}
\fig{calculo-I}{escalon}
\fig{calculo-I}{exponencial}
\par\bigskip
\fig{calculo-I}{extremo-derivada}
\fig{calculo-I}{extremos-globales-cuartica}
\fig{calculo-I}{extremos-globales-cubica}
\par\bigskip
\fig{calculo-I}{extremos-locales}
\fig{calculo-I}{factorizacion}
\fig{calculo-I}{funcion-acotada}
\par\bigskip
\fig{calculo-I}{funcion-constante}
\fig{calculo-I}{funcion-salto-punto}
\fig{calculo-I}{implicita-cubica}
\par\bigskip
\fig{calculo-I}{inflexion-cubo}
\fig{calculo-I}{inflexion-raiz-cubica}
\fig{calculo-I}{intervalos-finitos}
\par\bigskip
\fig{calculo-I}{intervalos-infinitos}
\fig{calculo-I}{inversas}
\fig{calculo-I}{limite-constante}
\par\bigskip
\fig{calculo-I}{limite-definicion}
\fig{calculo-I}{limite-identidad}
\fig{calculo-I}{logaritmo}
\par\bigskip
\fig{calculo-I}{monotonas}
\fig{calculo-I}{no-derivable-salto}
\fig{calculo-I}{no-derivable-varios}
\par\bigskip
\fig{calculo-I}{parabola-horizontal}
\fig{calculo-I}{paridad-impares}
\fig{calculo-I}{paridad-pares}
\par\bigskip
\fig{calculo-I}{pico-ejemplo}
\fig{calculo-I}{plano-cartesiano}
\fig{calculo-I}{postes}
\par\bigskip
\fig{calculo-I}{potencias}
\fig{calculo-I}{propiedad-supremo}
\fig{calculo-I}{quintica}
\par\bigskip
\fig{calculo-I}{racionales-irracionales}
\fig{calculo-I}{raices}
\fig{calculo-I}{raiz-mas-3}
\par\bigskip
\fig{calculo-I}{raiz-menos-3}
\fig{calculo-I}{raiz-reflejo-x}
\fig{calculo-I}{raiz-reflejo-y}
\par\bigskip
\fig{calculo-I}{raiz-x-mas-3}
\fig{calculo-I}{raiz-x-menos-3}
\fig{calculo-I}{reciproco}
\par\bigskip
\fig{calculo-I}{reciprocos}
\fig{calculo-I}{recta}
\fig{calculo-I}{recta-tangente}
\par\bigskip
\fig{calculo-I}{rectangulo-inscrito}
\fig{calculo-I}{restriccion-parabola}
\fig{calculo-I}{semicircunferencias}
\par\bigskip
\fig{calculo-I}{seno-extremos}
\fig{calculo-I}{seno-uno-entre-x}
\fig{calculo-I}{seno-x-tangente}
\par\bigskip
\fig{calculo-I}{sin-extremos}
\fig{calculo-I}{taylor-exponencial}
\fig{calculo-I}{transformacion-raiz}
\par\bigskip
\fig{calculo-I}{trigonometricas}
\fig{calculo-I}{valor-absoluto}
\fig{calculo-I}{valor-medio}
\par\bigskip
\fig{calculo-I}{valor-medio-generalizado}
\fig{calculo-I}{x-seno-uno-entre-x}
\fig{calculo-I}{x2-0-inf}
\par\bigskip
\fig{calculo-I}{x2-abierto-cerrado}
\fig{calculo-I}{x2-cerrado}
\fig{calculo-I}{x2-cerrado-abierto}
\par\bigskip
\fig{calculo-I}{x2-inf-0}
\fig{calculo-I}{x2-seno-uno-entre-x}
\fig{calculo-I}{x2-total}
\par\bigskip
\newpage

\section{Termodinámica}
{\small\color{black!60}45 figuras · \texttt{galeria/termodinamica/}}\par\bigskip

\fig{termodinamica}{PV-ejemplo}
\fig{termodinamica}{PV-trayectorias}
\fig{termodinamica}{S-U-dostercios}
\par\bigskip
\fig{termodinamica}{S-U-mitad}
\fig{termodinamica}{S-U-tercio}
\fig{termodinamica}{U-S-T}
\par\bigskip
\fig{termodinamica}{cajas-entropia}
\fig{termodinamica}{camino-PT}
\fig{termodinamica}{cavidad-resonante}
\par\bigskip
\fig{termodinamica}{cilindro-compuesto}
\fig{termodinamica}{cilindro-pistones}
\fig{termodinamica}{composicion}
\par\bigskip
\fig{termodinamica}{cuasiestatico}
\fig{termodinamica}{curva-Y}
\fig{termodinamica}{dipolo}
\par\bigskip
\fig{termodinamica}{esfera-proyeccion}
\fig{termodinamica}{extensivos-U}
\fig{termodinamica}{extensivos-V}
\par\bigskip
\fig{termodinamica}{familia-curvas}
\fig{termodinamica}{gas-recipiente}
\fig{termodinamica}{hipersuperficie-tangente}
\par\bigskip
\fig{termodinamica}{isotermas}
\fig{termodinamica}{ley-cero}
\fig{termodinamica}{pared-adiabatica}
\par\bigskip
\fig{termodinamica}{pared-diatermica-movil}
\fig{termodinamica}{pared-impermeable}
\fig{termodinamica}{pared-permeable}
\par\bigskip
\fig{termodinamica}{plano-S0}
\fig{termodinamica}{plano-U0}
\fig{termodinamica}{polimero}
\par\bigskip
\fig{termodinamica}{pozo-potencial}
\fig{termodinamica}{procesos-AB}
\fig{termodinamica}{region-TP}
\par\bigskip
\fig{termodinamica}{seis-subsistemas}
\fig{termodinamica}{sistema-compuesto}
\fig{termodinamica}{solenoide}
\par\bigskip
\fig{termodinamica}{sup-cuasiestatico}
\fig{termodinamica}{sup-equilibrio}
\fig{termodinamica}{sup-reversible}
\par\bigskip
\fig{termodinamica}{tangente-psi}
\fig{termodinamica}{tangentes}
\fig{termodinamica}{temperaturas-iguales}
\par\bigskip
\fig{termodinamica}{tipos-sistemas}
\fig{termodinamica}{trabajo-adiabatico}
\fig{termodinamica}{tres-cilindros}
\par\bigskip
\newpage

\section{Álgebra Lineal II}
{\small\color{black!60}32 figuras · \texttt{galeria/algebra-lineal-II/}}\par\bigskip

\fig{algebra-lineal-II}{6-3-minimos-cuadrados}
\fig{algebra-lineal-II}{6-5-elipse-ejes-rotados}
\fig{algebra-lineal-II}{6-6-proyeccion-ortogonal-vs-oblicua}
\par\bigskip
\fig{algebra-lineal-II}{6-7-circulo-elipse-valores-singulares}
\fig{algebra-lineal-II}{6-8-paraboloide-eliptico}
\fig{algebra-lineal-II}{accion-vector-propio}
\par\bigskip
\fig{algebra-lineal-II}{angulo-entre-vectores}
\fig{algebra-lineal-II}{cadena-ciudad-suburbios}
\fig{algebra-lineal-II}{diagrama-cociente-T-barra}
\par\bigskip
\fig{algebra-lineal-II}{diagrama-doble-dual}
\fig{algebra-lineal-II}{diagrama-phi-beta-LA}
\fig{algebra-lineal-II}{diagrama-valores-propios}
\par\bigskip
\fig{algebra-lineal-II}{discos-gerschgorin}
\fig{algebra-lineal-II}{distancia-punto-plano}
\fig{algebra-lineal-II}{elipse-cambio-coordenadas}
\par\bigskip
\fig{algebra-lineal-II}{generado-por-dos-vectores}
\fig{algebra-lineal-II}{gram-schmidt-paso}
\fig{algebra-lineal-II}{ley-paralelogramo}
\par\bigskip
\fig{algebra-lineal-II}{paralelepipedo-tres-vectores}
\fig{algebra-lineal-II}{paralelogramo-altura-comun}
\fig{algebra-lineal-II}{paralelogramo-diagonal-suma}
\par\bigskip
\fig{algebra-lineal-II}{paralelogramos-u-v}
\fig{algebra-lineal-II}{pendulo}
\fig{algebra-lineal-II}{plano-por-tres-puntos}
\par\bigskip
\fig{algebra-lineal-II}{recta-por-dos-puntos}
\fig{algebra-lineal-II}{reflexion-recta-y-2x}
\fig{algebra-lineal-II}{resorte-peso}
\par\bigskip
\fig{algebra-lineal-II}{rotacion-reflexion-proyeccion}
\fig{algebra-lineal-II}{sistemas-coordenados-relatividad}
\fig{algebra-lineal-II}{sistemas-mano-derecha-izquierda}
\par\bigskip
\fig{algebra-lineal-II}{suma-y-escalar-coordenadas}
\fig{algebra-lineal-II}{venn-bases}
\par\bigskip
\newpage

\section{Teoría analítica de números}
{\small\color{black!60}21 figuras · \texttt{galeria/teoria-numeros/}}\par\bigskip

\fig{teoria-numeros}{figura11-1}
\fig{teoria-numeros}{figura11-2}
\fig{teoria-numeros}{figura11-3}
\par\bigskip
\fig{teoria-numeros}{figura11-4}
\fig{teoria-numeros}{figura11-5}
\fig{teoria-numeros}{figura12-1}
\par\bigskip
\fig{teoria-numeros}{figura12-2}
\fig{teoria-numeros}{figura12-3}
\fig{teoria-numeros}{figura13-1}
\par\bigskip
\fig{teoria-numeros}{figura13-2}
\fig{teoria-numeros}{figura13-3}
\fig{teoria-numeros}{figura14-1}
\par\bigskip
\fig{teoria-numeros}{figura14-2}
\fig{teoria-numeros}{figura14-3}
\fig{teoria-numeros}{figura14-4}
\par\bigskip
\fig{teoria-numeros}{figura3-1}
\fig{teoria-numeros}{figura3-2}
\fig{teoria-numeros}{figura3-3}
\par\bigskip
\fig{teoria-numeros}{figura3-4}
\fig{teoria-numeros}{figura9-1}
\fig{teoria-numeros}{figura9-2}
\par\bigskip
\newpage

\section{Funciones trigonométricas e inversas}
{\small\color{black!60}12 figuras · \texttt{galeria/trigonometricas-inversas/}}\par\bigskip

\fig{trigonometricas-inversas}{arcocosecante}
\fig{trigonometricas-inversas}{arcocoseno}
\fig{trigonometricas-inversas}{arcocotangente}
\par\bigskip
\fig{trigonometricas-inversas}{arcosecante}
\fig{trigonometricas-inversas}{arcoseno}
\fig{trigonometricas-inversas}{arcotangente}
\par\bigskip
\fig{trigonometricas-inversas}{cosecante}
\fig{trigonometricas-inversas}{coseno}
\fig{trigonometricas-inversas}{cotangente}
\par\bigskip
\fig{trigonometricas-inversas}{secante}
\fig{trigonometricas-inversas}{seno}
\fig{trigonometricas-inversas}{tangente}
\par\bigskip
\newpage

\section{Lógica}
{\small\color{black!60}10 figuras · \texttt{galeria/logica/}}\par\bigskip

\fig{logica}{grafo-1}
\fig{logica}{grafo-2}
\fig{logica}{grafo-3}
\par\bigskip
\fig{logica}{grafo-4}
\fig{logica}{grafo-5}
\fig{logica}{grafos-encaje}
\par\bigskip
\fig{logica}{resolucion-cadena}
\fig{logica}{resolvente-definicion}
\fig{logica}{resolvente-ejemplo1}
\par\bigskip
\fig{logica}{resolvente-ejemplo2}
\par\bigskip
\newpage

\section{Álgebra Lineal I}
{\small\color{black!60}7 figuras · \texttt{galeria/algebra-lineal-I/}}\par\bigskip

\fig{algebra-lineal-I}{diagrama-cambio-base-Q}
\fig{algebra-lineal-I}{diagrama-cambio-base-ejemplo}
\fig{algebra-lineal-I}{diagrama-cambio-base-triangulo}
\par\bigskip
\fig{algebra-lineal-I}{diagrama-isomorfismo-Phi}
\fig{algebra-lineal-I}{diagrama-representacion-A}
\fig{algebra-lineal-I}{diagrama-representacion-T}
\par\bigskip
\fig{algebra-lineal-I}{diagrama-transformacion-dual}
\par\bigskip
\newpage

\section{Variable compleja}
{\small\color{black!60}7 figuras · \texttt{galeria/variable-compleja/}}\par\bigskip

\fig{variable-compleja}{argumento-w}
\fig{variable-compleja}{forma-polar}
\fig{variable-compleja}{plano-complejo}
\par\bigskip
\fig{variable-compleja}{raices-cubicas}
\fig{variable-compleja}{recta-vertical}
\fig{variable-compleja}{relaciones-disco-circulo}
\par\bigskip
\fig{variable-compleja}{suma-vectores}
\par\bigskip
\newpage

\end{document}
